\documentclass{situationist-preprint}

\title{Sartre's Paper-Knife}
\msid{EXI-09-b}
\msfamily{EXI-09}
\msdiscipline{Meaning, Death and the Absurd}
\mstype{Consequence}
\mskeywords{Sartre; artefacts; essence; real definition; function; defect}

\begin{document}
\maketitle

\begin{abstract}
Sartre illustrates the thesis that for artefacts essence precedes existence with a paper-knife, made by an artisan who had in mind both what it is for and how to make it. EXI-09-a argues that the thesis holds once \enquote{essence} is read as a defining project and artefact kinds are taken to be fixed by their makers' intentions. We draw three consequences from that result for the example itself. First, the paper-knife's real definition and its defining project coincide in content but differ in role, whereas a human being's real definition contains no defining project. Second, an object qualitatively identical to a paper-knife but formed without an intention is not a paper-knife. Third, only artefacts can be defective as what they are, in the sense in which a knife that does not cut is defective.
\end{abstract}

\section{The example}

Sartre's artisan makes the paper-knife from a concept, which includes both what the knife is for and a technique of production; the concept exists before the knife, and the knife exists because the concept was carried out \citep{Sartre1946}. EXI-09-a reads the first element of the concept, what the knife is for, as a defining project (Definition~2.1 of EXI-09-a), and derives from the theory of artefact kinds (Principle~P1 of EXI-09-a; \citealp{Hilpinen1992,Thomasson2003}) that for every artefact essence precedes existence (Lemma~3.1 of EXI-09-a). The technique of production plays no further role: it is how the intention is carried out, not what the artefact is for. We now draw out what the example shows once this reading is adopted.

\section{Three consequences}

\begin{proposition}[Two notions of essence]\label{p:two}
A paper-knife has a real definition, stating what kind of thing it is, and a defining project, stating what it is for, and the real definition includes the defining project. A human being has a real definition that does not include a defining project.
\end{proposition}

\begin{proof}
By Principle~P1 of EXI-09-a, a paper-knife is of its kind because it was made by someone intending to make a thing for cutting paper. So what it is, the content of its real definition in the sense of MET-13 \citep{Fine1994}, is a thing made to be for cutting paper, and the defining project is a constituent of it. For a human being, by the proof of Lemma~3.2 of EXI-09-a, the kind is fixed by biological origin and is compatible with indefinitely many lives; the real definition says what she is without saying what her life is for.
\end{proof}

Proposition~\ref{p:two} makes precise why Corollary~3.4 of EXI-09-a holds. MET-13 says that every object has a real definition; Sartre says that a human being has no essence prior to her existence. The two claims concern different things in the human case and coincide only in the artefact case, which is why Sartre's example can mislead: in the paper-knife, the two notions of essence cannot be told apart.

\begin{proposition}[The accidental knife]\label{p:swamp}
An object that is qualitatively identical to a paper-knife but was formed without any intention to make a thing for cutting paper is not a paper-knife, though it can be used as one.
\end{proposition}

\begin{proof}
By Principle~P1 of EXI-09-a, membership in an artefact kind requires that the object have been made by someone intending to make a thing of that kind. The accidental object was not. If someone later selects it and adopts it for cutting paper, a paper-knife may come into existence through that act, as in the driftwood case of EXI-09-a; but its essence then precedes its existence as a paper-knife, since the intention precedes it.
\end{proof}

\begin{proposition}[Defect]\label{p:defect}
A paper-knife that cannot cut paper is defective as a paper-knife. No human being is defective as a human being in virtue of what her life is for.
\end{proposition}

\begin{proof}
By Lemma~3.1 of EXI-09-a, the paper-knife's defining project is fixed before it exists, so there is a standard of success as what it is that does not depend on how it is used; a knife that cannot cut fails it. By Lemma~3.2 of EXI-09-a, no defining project for a human life is fixed before it exists; whatever standard of success her life has as hers is fixed by the ends she adopts (Principle~P2 of EXI-09-a). A life cannot fail as a human life by departing from a project it never had.
\end{proof}

Proposition~\ref{p:defect} does not say that lives cannot be evaluated. A life can be imprudent, wicked or, in the sense of EXI-01, less meaningful than it might have been, because it engages little with what is worthwhile. These evaluations apply standards of value, not a defining project, and they do not make a life a defective specimen of anything.

\begin{objection}[Natural functions]
Organs have functions without makers: the heart is for pumping blood, and a heart that cannot pump is defective. So defect does not require a prior intention, and a human being might be defective as an organism.
\end{objection}

\begin{reply}
We grant that biological function grounds a standard for organs and perhaps for organisms. But a defining project, in Definition~2.1 of EXI-09-a, is a determination of what a thing is for in the light of which its activities count as succeeding or failing as what it is; the paper-knife's project is fixed by intention and is not an evolutionary history. A human being may be unhealthy as an organism; that is not a failure of her life to be what it is for. Proposition~\ref{p:defect} concerns the second.
\end{reply}

\begin{objection}[Mass production]
Modern paper-knives are not made by an artisan with a concept but by machines running a design. No one intends this knife.
\end{objection}

\begin{reply}
The intention of Principle~P1 of EXI-09-a can be general: the designer and the manufacturer intend that each item produced by the process be a paper-knife, and each is made in carrying out that intention. Precedence in Definition~2.2 of EXI-09-a requires that a defining project be fixed before the object exists by a determination on which its existence depends, not that the determination mention it singly.
\end{reply}

\section{Conclusion}

The paper-knife is a good example of the precedence of essence because in an artefact the real definition and the defining project coincide. Its value for Sartre lies in the contrast: a human being has a real definition but no prior defining project, and so cannot be defective in the way the knife can.

\begin{thebibliography}{9}\small
\bibitem[Fine(1994)]{Fine1994} Fine, K. (1994). Essence and modality. \emph{Philosophical Perspectives}, 8, 1--16.
\bibitem[Hilpinen(1992)]{Hilpinen1992} Hilpinen, R. (1992). Artifacts and works of art. \emph{Theoria}, 58(1), 58--82.
\bibitem[Sartre(1946)]{Sartre1946} Sartre, J.-P. (1946). \emph{L'existentialisme est un humanisme}. Nagel.
\bibitem[Thomasson(2003)]{Thomasson2003} Thomasson, A.~L. (2003). Realism and human kinds. \emph{Philosophy and Phenomenological Research}, 67(3), 580--609.
\end{thebibliography}

\end{document}
